% Source: https://artofproblemsolving.com/community/c3h1975278_descartes_rule_of_signs_clarification
% Stable identity: aops-topic-1975278-q01
Since $p(1)=0$, divide by $x-1$: $x^3-3x+2=(x-1)(x^2+x-2)=(x-1)^2(x+2)$. So $p$ has the double root $1$ and the simple root $-2$.

The coefficients of $p(x)$ are $+1,-3,+2$ (the zero $x^2$ coefficient is skipped), with $2$ sign changes. The positive roots counted with multiplicity are $1,1$: two of them. Without counting multiplicity there would be only one positive root, and $2-1$ is odd, which would contradict the rule; so the rule must be read with multiplicity.

For negative roots, $p(-x)=-x^3+3x+2$ has coefficients $-1,+3,+2$ with $1$ sign change, and $p$ has exactly one negative root, $-2$.

So the counts agree exactly: $\boxed{2\text{ positive roots (the double root }1)\text{ and }1\text{ negative root }(-2)}$, matching $2$ and $1$ sign changes.
